\centerline{\Huge \bf Начинаем синусовать}
\centerline{\Huge \bf }

\begin{flushright}
        \scriptsize{Васильков. Всей душой рад служить вам;\\ но чему же я могу учить Лидию Юрьевну? Сферической тригонометрии?\\Островский. <<Бешеные деньги>> 1870}
\end{flushright}


\problem{1}
 В треугольнике провели (а) высоты (б) биссектрисы. Выразите длины \textbf{ВСЕХ} получившихся отрезков через $\angle A = \alpha, \angle B = \beta, \angle C = \gamma$ и радиус описанной окружности $R$.

\problem{2}
 \textbf{Важная лемма.} В треугольнике $ABC$ на прямой $AC$ выбрана точка $D$. Докажите, что $\frac{\sin\angle ABD}{\sin \angle DBC} = \frac{AD \cdot BC}{DC \cdot AB}$

\problem{3}
В остроугольном треугольнике $ABC$ проведены высоты $AA_1$ и $BB_1$. Точка $O$ --- центр описанной окружности. Докажите, что расстояние от точки $A_1$ до прямой $BO$ такое же как расстояние от точки $B_1$ до прямой $AO$.

\problem{4}
Пусть $ABCD$ --- вписанный четырехугольник, $M, N$ --- середины дуг $AB$ и $CD$. Точки $I, J$ --- центры вписанный окружностей $ABD$ и $CBD$. Докажите, что $MN$ делит $IJ$ пополам.

\problem{5}
\textbf{Теорема Наполеона.} На сторонах треугольника $ABC$ наружу построены равносторонние треугольники. Докажите, что их центр образуют правильный треугольник. 


\begin{flushright}
	\textit{Примечание: поворотом пользоваться нельзя}	
\end{flushright}

\problem{6}
Точка $X$ лежит внутри правильного треугольника $ABC$. Точки $A_1, B_1, C_1$ симметричны $X$ относительно сторон $BC, AC, AB$ соответственно. Докажите, что прямые $AA_1, BB_1, CC_1$ пересекаются в одной точке.

